% =====================================================================
%  resumo.tex  ->  Quadro-resumo executivo + Resumo + Palavras-chave
%  Texto identico ao .docx. O quadro-resumo reaproveita frases verbatim
%  da Conclusao (nenhum conteudo novo foi criado).
% =====================================================================

\begin{quadroresumo}[Em resumo]
O Backtracking é uma estratégia baseada em tentativa e erro inteligente. Ao invés de explorar todas as possibilidades cegamente, o algoritmo constrói soluções passo a passo, verificando constantemente se ainda existe chance de alcançar o resultado desejado. Sempre que identifica um caminho inviável, retorna para a última decisão válida e tenta uma alternativa diferente.
\end{quadroresumo}

\subsection*{Resumo}

Diversos problemas computacionais exigem a exploração de múltiplas possibilidades até que uma solução válida seja encontrada. Em muitos casos, testar todas as alternativas é inviável devido à grande quantidade de combinações possíveis. Para lidar com esse desafio, foi desenvolvida a técnica de Backtracking, um método que constrói soluções gradualmente e retorna para estados anteriores sempre que identifica que um determinado caminho não levará ao resultado desejado.

Este trabalho apresenta os conceitos fundamentais do Backtracking, suas características, seu funcionamento e aplicações práticas, utilizando exemplos simples e intuitivos para facilitar o entendimento do leitor.

\subsection*{Palavras-chave}

Backtracking, Grafos, Algoritmos, Busca Exaustiva, Árvores de Busca.
